\section{当前实现的问题分析}
\label{sec:h6-analysis}

本节分析第~\ref{sec:h6-baseline} 节所述的旧版 H6 基线。

\begin{enumerate}[itemsep=0.55em]
  \item \textbf{83 维输入中有 3 维可重建。}
  由式~\eqref{eq:redundant}，常规控制帧中，最新观测的前 3 维等于 \code{cmd_ori} 中两组三维向量的和。保留原始指令和上一帧修正，即可从 MLP 输入中省去合成指令三维。去冗余对训练性能的影响仍需消融验证。\ev{C2,C4}

  \item \textbf{仿真施力与头部附近牵引存在差异。}
  旧版仿真在质心施力。实际从头部附近牵引时，拉点相对质心的偏移会产生等效力矩，见式~\eqref{eq:force-transform}。此外，参考映射使用脉冲峰值，机器人只在控制周期的部分时间内受力。\ev{C7}

  \item \textbf{平移参考修正越过阈值就饱和。}
  式~\eqref{eq:force-map} 中，任一平移轴的力幅值刚超过 $17.5\,\mathrm{N}$，限幅前的修正幅值就已超过 $1.75\,\mathrm{m/s}$，直接达到该轴的限幅。例如，$18\,\mathrm{N}$ 和 $50\,\mathrm{N}$ 的前向力都对应 $1.2\,\mathrm{m/s}$ 的参考修正，阈值以上的力幅值差异无法反映到平移参考中。\ev{C3,C5}

  \item \textbf{运动跟踪奖励的设计与牵引任务不匹配。}
  底层负责跟踪给定速度，牵引 MLP 从观测中学习与外力对应的速度修正。当前平移和偏航跟踪奖励仍计入牵引 PPO 回报，而底层已冻结。其跟踪目标为原始指令加上一帧 MLP 修正，策略可能通过选取底层更容易执行的修正来提高跟踪分，即使该修正偏离外力参考。这两项跟踪奖励应留在底层控制训练中。\ev{C6,C8}

  \item \textbf{撤力后的实际减速过程与停车判据未明确。}
  旧版基线在撤力后将参考修正清零。纯牵引部署时外部原始指令为零，停车取决于网络输出是否回零，以及底座如何减速。第~\ref{sec:evaluation} 节依据撤力后的位移、航向、直立状态和底座执行指令评估停稳。\ev{C5,C6,C7,C11}

  \item \textbf{力标签比实际施力提前一帧。}
  当前控制步先用旧力完成仿真，再更新力标签，并用新标签计算本帧的 \code{compliance}。例如，撤力边界这一帧仍受最后一次非零力脉冲，奖励目标却已归零，到下一帧才停止施力。这是切换时相差一个控制步（$0.02\,\mathrm{s}$）的时序错误，基线因复现原实现而保留了这一行为。\ev{C6,C7,C8}
\end{enumerate}
